% Options for packages loaded elsewhere
\PassOptionsToPackage{unicode}{hyperref}
\PassOptionsToPackage{hyphens}{url}
\documentclass[
  ukrainian,
]{article}
\usepackage{xcolor}
\usepackage[margin=2.5cm]{geometry}
\usepackage{amsmath,amssymb}
\setcounter{secnumdepth}{-\maxdimen} % remove section numbering
\usepackage{iftex}
\ifPDFTeX
  \usepackage[T1]{fontenc}
  \usepackage[utf8]{inputenc}
  \usepackage{textcomp} % provide euro and other symbols
\else % if luatex or xetex
  \usepackage{unicode-math} % this also loads fontspec
  \defaultfontfeatures{Scale=MatchLowercase}
  \defaultfontfeatures[\rmfamily]{Ligatures=TeX,Scale=1}
\fi
\usepackage{lmodern}
\ifPDFTeX\else
\setmainfont{DejaVu Serif}
\setsansfont{DejaVu Sans}
\setmonofont{DejaVu Sans Mono}
\setmathfont{STIX Two Math}
\fi
\ifPDFTeX\else
  % xetex/luatex font selection
\fi
% Use upquote if available, for straight quotes in verbatim environments
\IfFileExists{upquote.sty}{\usepackage{upquote}}{}
\IfFileExists{microtype.sty}{% use microtype if available
  \usepackage[]{microtype}
  \UseMicrotypeSet[protrusion]{basicmath} % disable protrusion for tt fonts
}{}
\makeatletter
\@ifundefined{KOMAClassName}{% if non-KOMA class
  \IfFileExists{parskip.sty}{%
    \usepackage{parskip}
  }{% else
    \setlength{\parindent}{0pt}
    \setlength{\parskip}{6pt plus 2pt minus 1pt}}
}{% if KOMA class
  \KOMAoptions{parskip=half}}
\makeatother
\usepackage{longtable,booktabs,array}
\newcounter{none} % for unnumbered tables
\usepackage{calc} % for calculating minipage widths
% Correct order of tables after \paragraph or \subparagraph
\usepackage{etoolbox}
\makeatletter
\patchcmd\longtable{\par}{\if@noskipsec\mbox{}\fi\par}{}{}
\makeatother
% Allow footnotes in longtable head/foot
\IfFileExists{footnotehyper.sty}{\usepackage{footnotehyper}}{\usepackage{footnote}}
\makesavenoteenv{longtable}
\usepackage{graphicx}
\graphicspath{{/Users/kalinicenkovitalijmikolajovic/.codex/worktrees/5c15/Script_for_translate_book/Output_Final/}}
\setkeys{Gin}{width=0.82\linewidth,height=0.82\textheight,keepaspectratio}
\makeatletter
\newsavebox\pandoc@box
\newcommand*\pandocbounded[1]{% scales image to fit in text height/width
  \sbox\pandoc@box{#1}%
  \Gscale@div\@tempa{\textheight}{\dimexpr\ht\pandoc@box+\dp\pandoc@box\relax}%
  \Gscale@div\@tempb{\linewidth}{\wd\pandoc@box}%
  \ifdim\@tempb\p@<\@tempa\p@\let\@tempa\@tempb\fi% select the smaller of both
  \ifdim\@tempa\p@<\p@\scalebox{\@tempa}{\usebox\pandoc@box}%
  \else\usebox{\pandoc@box}%
  \fi%
}
% Set default figure placement to htbp
\def\fps@figure{htbp}
\makeatother
\ifLuaTeX
\usepackage[bidi=basic,shorthands=off,provide=*]{babel}
\else
\usepackage[bidi=default,shorthands=off,provide=*]{babel}
\fi
\ifLuaTeX
  \usepackage{selnolig} % disable illegal ligatures
\fi
\setlength{\emergencystretch}{3em} % prevent overfull lines
\providecommand{\tightlist}{%
  \setlength{\itemsep}{0pt}\setlength{\parskip}{0pt}}
\usepackage{bookmark}
\IfFileExists{xurl.sty}{\usepackage{xurl}}{} % add URL line breaks if available
\urlstyle{same}
\hypersetup{
  pdflang={uk},
  hidelinks,
  pdfcreator={LaTeX via pandoc}}

\author{}
\date{}

\begin{document}

\section{Задача теплопровідності з похідною в граничній
умові}\label{ux437ux430ux434ux430ux447ux430-ux442ux435ux43fux43bux43eux43fux440ux43eux432ux456ux434ux43dux43eux441ux442ux456-ux437-ux43fux43eux445ux456ux434ux43dux43eux44e-ux432-ux433ux440ux430ux43dux438ux447ux43dux456ux439-ux443ux43cux43eux432ux456}

Розглянемо пристрій, показаний на рис. 7.1. Припустимо, що верхній
кінець стрижня підтримується при нульовій температурі, тобто \(u(0, t)\)
= 0, а нижній кінець занурений у воду, температура якої також стала і
дорівнює нулю (у цьому випадку --- нуль Р и с

\pandocbounded{\includegraphics[keepaspectratio]{images/_page_50_Picture_4.jpeg}}

. 7.1. Схематичне представлення зсувної задачі.

\pandocbounded{\includegraphics[keepaspectratio]{images/_page_50_Figure_6.jpeg}}

Рис. 7.2. Природа кривих, які задовольняють

\[
(\Gamma Y) =
\begin{cases} u(0, t) = 0, \\ u_x(1, t) = -hu(1, t). \end{cases}
\]

відповідає певній еталонній температурі). Теплообмін між стрижнем і
водою відбувається згідно із законом Ньютона, який записується у x=1 як
наступна гранична умова:

\[
u_x(1, t) = -hu(1, t).
\]

Тепер припустимо, що початкова температура стрижня дорівнює \(u(x, 0)\)
= x, але одразу при t \textgreater{} 0 включаємо граничні умови. Щоб
визначити температуру в стрижні, необхідно розв'язати проблему зміщення

\begin{verbatim}
\begin{aligned}
(\text{УЧП})\quad & u_t = \alpha^2 u_{xx}, && 0 < x < 1,\ 0 < t < \infty, \\
(7.1)\ (\text{ГУ})\quad & u(0,t) = 0,\quad u_x(1,t) + h u(1,t) = 0, && 0 < t < \infty, \\
(\text{НУ})\quad & u(x,0) = x, && 0 \leq x \leq 1.
\end{aligned}
\end{verbatim}

Щоб використати метод розділення змінних, розбиємо процес розв'язання на
кілька етапів.

\textbf{КРОК} 1. (Перетворення диференціального рівняння в похідних у
два звичайних диференціальних рівняння)

\begin{verbatim}
Підставляючи u(x, t) = X(x)T(t) у рівняння (7.1), отримуємо XT' = \alpha^2 X''T Поділивши інші частини цього рівняння на \alpha^2 XT отримуємо
\end{verbatim}

\[
\frac{T'}{\alpha^2 T} = \frac{X''}{X} .
\]

\begin{verbatim}
Оскільки ліва частина останнього рівняння залежить лише від часу, а права — лише від x (і оскільки x і t незалежні), обидві сторони рівняння мають бути константами. Прирівнюючи обидві частини до однієї сталої \mu , отримуємо
\end{verbatim}

(7.2)

\[
T' - \mu \alpha^{2} T = 0, X'' - \mu X = 0.
\]

Процес розщеплення змінних завершено.

КРОК 2. (Визначення константи відокремлення)

\begin{verbatim}
По-перше, константа і не повинна бути додатною, інакше ^{1} ) функція T\left( t\right) буде зростати експоненційно, що призведе до нескінченного підвищення температури u=XT. Очевидно, що таке рішення не має фізичного значення.
\end{verbatim}

\begin{verbatim}
По-друге, припустимо, що \mu = 0 . Тоді
\end{verbatim}

\[
X'' = 0
\]

і, отже.

\[
X(x) = A + \hat{B}x.
\]

Однак, оскільки граничні умови вимагають, що

\[
u(0, t) = X(0) T(t) = 0,
\]

\(u(1, t)\) = (hX(1) + X'(1)) \(T(t)\) = 0, Ми одразу дійшли висновку,
що

\[
X(0) = 0
\Rightarrow A = 0 X'(1) + hX(1) = 0 \Rightarrow B = 0 .
\]

\begin{verbatim}
Це означає, що u(x, t) \equiv 0 . Іншими словами, лише нульове розв'язання відповідає нульовій константі розділення. Ми відкинемо його, оскільки нас цікавлять ненульові розв'язки.
\end{verbatim}

\begin{verbatim}
Якщо, нарешті, \mu < 0 , можна поставити \mu = -\lambda^2 і переписати рівняння (7.2) у вигляді
\end{verbatim}

\[
T' + \lambda^2 \alpha^2 T = 0,
X' + \lambda^2 X = 0.
\]

1) Розв'язок \(X(x)\) відповідної задачі з граничним значенням буде лише
тривіальним, \(X(x)\) = 0. --- Ед. Ед.

Загальні розв'язки цих рівнянь легко знайти:

\[
T(t) = Ae^{-(\lambda \alpha)^2 t},
X(x) = B\sin(\lambda x) + C\cos(\lambda x).
\]

Отже, будь-яка функція провідника

(7.3)

\[
u(x, t) = e^{-(\lambda \alpha)^2 t} [A \sin(\lambda x) + B \cos(\lambda x)]
\]

\begin{verbatim}
при будь-яких значеннях \lambda A і B задовольняє рівняння (читач може перевірити це самостійно простою заміною).
\end{verbatim}

Тепер нам потрібно з'ясувати, які з цих функцій задовольняють граничним
умовам

(7.4)

\[
u(0, t) = 0, u_x(1, t) + hu(1, t) = 0.
\]

\begin{verbatim}
Підстановка розв'язків виду (7.3) у граничні умови (7.4) дозволяє знайти відношення, яким мають задовольняти значення \lambda , \Lambda та B:
\end{verbatim}

\[
Be^{-(\lambda\alpha)^2 t} = 0 \Rightarrow B = 0,
\]

\begin{verbatim}
A\lambda e^{-(\lambda\alpha)^2 t} \cos \lambda + hAe^{-(\lambda\alpha)^2 t} \sin \lambda = 0. Після простих алгебранхічних перетворень останнє рівняння зводиться до вигляду
\end{verbatim}

\[
tg \lambda = -\lambda/\hbar
\]

.

\begin{verbatim}
Іншими словами, щоб визначити \lambda , потрібно знайти точки перетину графів функцій tg \lambda і — \lambda/h (рисунок 7.3).
\end{verbatim}

\begin{verbatim}
Числові значення значень \lambda_1 , \lambda_2 , ... можна знайти за допомогою комп'ютера, якщо задано значення h Р и с
\end{verbatim}

\pandocbounded{\includegraphics[keepaspectratio]{images/_page_52_Figure_15.jpeg}}

\begin{verbatim}
. 7.3. Перетин графів функцій tg(\lambda) і -\lambda/h .
\end{verbatim}

викликати власні значення задачі граничного значення

(7.5)

\[
X'' + \lambda^{2}X = 0,
\begin{cases} X(0) = 0 \\ X'(1) + hX(1) = 0. \end{cases}
\]

Іншими словами, з цими значеннями параметра \(\lambda\) існують
нетривіальні розв'язки задачі (7.5). Власні значення задачі (7.5)
визначаються як корені рівняння \(\log \lambda = -\lambda/h\) . Їхні
числові значення знаходилися на комп'ютері для випадку h = 1. Таблиця
7.1 показує результати розрахунків перших кількох власних значень.

{\def\LTcaptype{none} % do not increment counter
\begin{longtable}[]{@{}llll@{}}
\toprule\noalign{}
n & \(\lambda_B\) & rt . & \(\lambda_n\) \\
\midrule\noalign{}
\endhead
\bottomrule\noalign{}
\endlastfoot
123 & 2,024,917,98 & 45 & 11,0814,20 \\
\end{longtable}
}

Розв'язки задачі (7.5), що відповідають власним значенням \(\lambda_n\)
, називаються власними функціями \(X_n(x)\) . Для розглянутого завдання
власні функції виглядають так:

\[
X_n(x) = \sin(\lambda_n x)
\]

і зображені на рис. 7.4.

\pandocbounded{\includegraphics[keepaspectratio]{images/_page_53_Figure_10.jpeg}}

Рис. 7.4. Власні функції \(X_n(x)\) завдання (7.5) при h=1; Усі власні
функції задовольняють граничні умови.

КРОК 3, (Пошук фундаментальних рішень)

Тепер у нас є нескінченний набір функцій (фундаментальних розв'язків) у
нашому розпорядженні

\[
u_n(x, t) = X_n(x) T_n(t) = e^{-(\lambda_n x)^2 t} \sin(\lambda_n x),
\]

кожна з яких задовольняє рівняння та граничні умови. Залишився ще один
крок --- загалом

\[
u(x, t) = \sum_{n=1}^{\infty} a_n X_n(x) T_n(t) =
= \sum_{n=1}^{\infty} a_n e^{-(\lambda_n c t^2 t)} \sin(\lambda_n x)
\]

Вибрати коефіцієнти \(a_n\) так, щоб функція \(u(x, t)\) задовольняла
початкову умову \(u(x, 0)\) = x (у цьому випадку функція \(u(x, t)\)
задовольнятиме рівняння та граничні умови, оскільки, по-перше, кожен
член суми їх задовольняє, а по-друге, і рівняння, і граничні умови є
лінійними та однорідними).

Отже, ймовірність \(a_n\) слід визначати з цієї умови

(7 6)

\[
u(x, 0) = x = \sum_{n=1}^{\infty} a_n \sin(\lambda_n x).
\]

Ми маємо зробити останній крок.

КРОК 4. (Розклад початкової умови на ряд за власними функціями)

\begin{verbatim}
Щоб знайти коефіцієнти a_n у розкладі (7,6) на власні функції, помножте обидві частини цього співвідношення на \sin(\lambda_m x) і інтегруйте від 0 до 1. У результаті ми отримуємо
\end{verbatim}

\[
\int_{0}^{2} \xi \sin(\lambda_{m} \xi) d\xi = \sum_{n=1}^{\infty} a_{n} \int_{0}^{1} \sin(\lambda_{n} \xi) \sin(\lambda_{m} \xi) d\xi =
= a_{m} \int_{0}^{1} \sin^{2}(\lambda_{m} \xi) d\xi =
= a_{m} \left( \frac{\lambda_{m} - \sin \lambda_{m} \cos \lambda_{m}}{2\lambda_{m}} \right).
\]

Розв'язавши останнє відношення відносно \(a_m\) і замінивши m на n,
отримуємо бажаний результат:

(7.7)

\[
a_n = \frac{2\lambda_n}{(\lambda_n - \sin \lambda_n \cos \lambda_n)} \int_0^{\infty} \xi \sin (\lambda_n \xi) d\xi.
\]

Таблиця 7.2.

\section{\texorpdfstring{Шанси \(a_n\) у формулі
(7.8)}{Шанси a\_n у формулі (7.8)}}\label{ux448ux430ux43dux441ux438-a_n-ux443-ux444ux43eux440ux43cux443ux43bux456-7.8}

n \textbar{} an \textbar{} n \textbar{} a\_n \textbar{}

\textbar-------------\textbar-----------------------\textbar---------------\textbar----------------\textbar{}

123 \textbar{} 0,240,22-0,03 \textbar{} 45 \textbar{} -0.11-0.09
\textbar{}

Отже, розв'язок задачі (7.1) записується у вигляді

(7.8)

\[
u(x, t) = \sum_{n=1}^{\infty} a_n e^{-i\lambda_n \alpha_n^{2} t} \sin(\lambda_n x),
\]

де коефіцієнти \(a_n\) задаються формулами (7.7). Перші п'ять
коефіцієнтів \(a_n\) нашої задачі були знайдені чисельно. Їхні значення
наведені в Таблиці 1. 7.2.

\pandocbounded{\includegraphics[keepaspectratio]{images/_page_55_Figure_8.jpeg}}

Рис. 7.5. Графічне представлення функції (7.8); межа умовно
задовольняється у всі моменти часу.

Відповідно, перші три члени розв'язання змішаної задачі

(7.9)

\begin{verbatim}
\begin{aligned} u_t &= u_{xx}, & 0 < x < 1, & 0 < t < \infty, \\ u_t &= 0, & 0 < t < \infty, \\ u_x &= 0, & 0 < t < \infty, \\ u_x &= 0, & 0 < t < \infty, \\ u_x &= 0, & 0 < t < \infty, \\ u_x &= 0, & 0 < t < \infty, \\ u_x &= 0, & 0 < t < \infty, \\ u_x &= 0, & 0 < t < \infty, \\ u_x &= 0, & 0 < t < \infty, \\ u_x &= 0, & 0 < t < \infty, \\ u_x &= 0, & 0 < t < \infty, \\ u_x &= 0, & 0 < t < \infty, \\ u_x &= 0, & 0 < t < \infty, \\ u_x &= 0, & 0 < t < \infty, \\ u_x &= 0, & 0 < t < \infty, \\ u_x &= 0, & 0 < t < \infty, \\ u_x &= 0, & 0 < t < \infty, \\ u_x &= 0, & 0 < t < \infty, \\ u_x &= 0, & 0 < t < \infty, \\ u_x &= 0, & 0 < t < \infty, \\ u_x &= 0, & 0 < t < \infty, \\ u_x &= 0, & 0 < t < \infty, \\ u_x &= 0, & 0 < t < \infty, \\ u_x &= 0, & 0 < t < \infty, \\ u_x &= 0, & 0 < t < \infty, \\ u_x &= 0, & 0 < t < \infty, \\ u_x &= 0, & 0 < t < \infty, \\ u_x &= 0, & 0 < t < \infty, \\ u_x &= 0, & 0 < t < \infty, \\ u_x &= 0, & 0 < t < \infty, \\ u_x &= 0, & 0 < t < \infty, \\ u_x &= 0, & 0 < t < \infty, \\ u_x &= 0, & 0 < t < \infty, \\ u_x &= 0, & 0 < t < \infty, \\ u_x &= 0, & 0 < t < \infty, \\ u_x &= 0, & 0 < t < \infty, \\ u_x &= 0, & 0 < t < \infty, \\ u_x &= 0, & 0 < t < \infty, \\ u_x &= 0, & 0 < t < \infty, \\ u_x &= 0, & 0 < t < \infty, \\ u_x &= 0, & 0 < t < \infty, \\ u_x &= 0, & 0 < t < \infty, \\ u_x &= 0, & 0 < t < \infty, \\ u_x &= 0, & 0 < t < \infty, \\ u_x &= 0, & 0 < t < \infty, \\ u_x &= 0, & 0 < t < \infty, \\ u_x &= 0, & 0 < t < \infty, \\ u_x &= 0, & 0 < t < \infty, \\ u_x &= 0, & 0 < t < \infty, \\ u_x &= 0, & 0 < t < \infty, \\ u_x &= 0, & 0 < t < \infty, \\ u_x &= 0, & 0 < t < \infty, \\ u_x &= 0, & 0 < t < \infty, \\ u_x &= 0, & 0 < t < \infty, \\ u_x &= 0, & 0 < t < \infty, \\ u_x &= 0, & 0 < t < \infty, \\ u_x &= 0, & 0 < t < \infty, \\ u_x &= 0, & 0 < t < \infty, \\ u_x &= 0, & 0 < t < \infty, \\ u_x &= 0, & 0 < t < \infty, \\ u_x &= 0, & 0 < t < \infty, \\ u_x &= 0, & 0 < t < \infty, \\ u_x &= 0, & 0 < t < \infty, \\ u_x &= 0, & 0 < t < \infty, \\ u_x &= 0, & 0 < t < \infty, \\ u_x &= 0, & 0 < t < \infty, \\ u_x &= 0, & 0 < t < \infty, \\ u_x &= 0, & 0 < t < \infty, \\ u_x &= 0, & 0 < t < \infty, \\ u_x &= 0, & 0 < t < \infty, \\ u_x &= 0, & 0 < t < \infty, \\ u_x &= 0, & 0 < t < \infty, \\ u_x &= 0, & 0 < t < \infty, \\ u_x &= 0, & 0 < t < \infty, \\ u_x &= 0, & 0 < t < \infty, \\ u_x &= 0, & 0 < t < \infty, \\ u_x &= 0, & 0 < t < \infty, \\ u_x &= 0, & 0 < t < \infty, \\ u_x &= 0, & 0 < t < \infty, \\ u_x &= 0, & 0 < t < \infty, \\ u_x &= 0, & 0 < t
\end{verbatim}

мають вигляд

\[
u(x, t) = = 0.24e^{-4t} \sin(2.02x) + 0.22e^{-24t} \sin(4.9x) - 0.03e^{-48.3t} \sin(7.98x) + \dots
\]

Фіг. На рисунку 7-5 показані графіки цього розв'язку для кількох точок
часу. Рекомендуємо читачеві відповісти на такі питання: Чи відповідає
тип рішення, яке ви отримуєте, з вашими інтуїтивними ідеями? Чи
виконуються граничні умови?

\section{ЗАУВАЖЕННЯ}\label{ux437ux430ux443ux432ux430ux436ux435ux43dux43dux44f}

Задача власних значень (7.5) є особливим випадком більш загальної
задачі, яку зазвичай називають проблемою Штурма-Ліувіля:

(7.10)

\[
\begin{aligned} & (\text{ОДУ}) \quad [p(x)y']' - q(x)y + \lambda r(x)y = 0, \quad 0 < x < 1, \\ & (\alpha_1 y(0) + \beta_1 y'(0) = 0, \\ & (\alpha_2 y(1) + \beta_1 y'(1) = 0. \end{aligned}
\]

Якщо розв'язати задачу з однорідними граничними умовами методом
розділення змінних, то функції \(X(x)\) завжди знаходяться як розв'язки
деякої задачі Штурма--Ліувіля. Очевидно, що задача власних значень (7.5)
є особливим випадком задачі (7.10).

Перелічимо основні властивості задачі (7.10), які були встановлені
Штурмом і Ліувілем за певних природних обмежень.

функції \(p(x)\), \(q(x)\) та \(r(x)\).

\begin{enumerate}
\def\labelenumi{\arabic{enumi}.}
\tightlist
\item
  Існує нескінченна послідовність власних значень, що задовольняють
  нерівності
\end{enumerate}

\[
\lambda_1 < \lambda_2 < \lambda_3 < \ldots < \lambda_n < \ldots \to \infty
\]

.

\begin{verbatim}
2. Кожне власне значення \lambda_n відповідає одній (точній до константного множника) ненульовій власній функції y_n(x) .
\end{verbatim}

\begin{verbatim}
3. Якщо y_n(x) і y_m(x) дві різні власні функції, що відповідають різним власним значенням \lambda_n і \lambda_m , то вони є ортогональними з вагою r(x) на сегменті [0, 1], тобто вони задовольняють це відношення
\end{verbatim}

\[
\int_{0}^{1} r(x) y_{n}(x) y_{m}(x) dx = 0.
\]

Більше деталей про проблему Штурма-Ліувіля можна знайти в літературі,
наведеній наприкінці книги.

1) Див., наприклад, Тіхоног А. Н., Самарський А. А. Рівняння
математичної фізики. --- Москва: Наука, 1972. --- ред. Ед.

\section{ЗАВДАННЯ}\label{ux437ux430ux432ux434ux430ux43dux43dux44f}

\begin{enumerate}
\def\labelenumi{\arabic{enumi}.}
\tightlist
\item
  Розв'язати наступну задачу для рівняння теплопровідності методом
  розкладу
\end{enumerate}

(PPP)

\[
u_1 = u_{xx}, \quad 0 < x < 1, \quad 0 < t < \infty,
\]

(GI)

\[
\begin{cases} u(0, t) = 0, \\ u_x(1, t) = 0, \\ u(x, 0) = x, \quad 0 \le x \le 1. \end{cases}
\]

Чи відповідає рішення, яке ви знаходите, з інтонативними
представленнями? Що таке стаціонарне рішення?

\begin{enumerate}
\def\labelenumi{\arabic{enumi}.}
\setcounter{enumi}{1}
\tightlist
\item
  Знайдіть власні цінності та функції проблеми Ассолта-Ліувіля
\end{enumerate}

(ODY)

\[
X' + \lambda X = 0, \quad 0 < x < 1,
\]

(FY)

\[
\begin{cases} X(0) = 0, \\ X'(1) = 0. \end{cases}
\]

Задайте специфічні вирази для функцій \(p(x)\), \(q(x)\) та \(r(x)\) із
загальної задачі Штурма--Ліувіля для цього випадку.

\begin{enumerate}
\def\labelenumi{\arabic{enumi}.}
\setcounter{enumi}{2}
\tightlist
\item
  Розв'язати задачу теплопровідності для планки з термоізольованими
  межами
\end{enumerate}

(PPP)

\[
u_t = u_{xx}, \quad 0 < x < 1, \quad 0 < t < \infty,
\]

(GI)

\[
\begin{cases} u_x (0, t) = 0, \\ u_x (1, t) = 0, \end{cases} \quad 0 < t < \infty,
\]

\begin{verbatim}
(ГУ) u(x, 0) = x, \quad 0 \le x \le 1. Чи відповідає рішення, яке ви отримуєте, з вашою інтерпретацією цієї проблеми? Що таке стаціонарне рішення? Яке його фізичне значення?
\end{verbatim}

\begin{enumerate}
\def\labelenumi{\arabic{enumi}.}
\setcounter{enumi}{3}
\tightlist
\item
  Знайдіть власні цінності та функції завдань
\end{enumerate}

(ОДА)

\[
X'' + \lambda X = 0, \quad 0 < x < 1,
\]

(GI)

\[
\begin{cases} X'(0) = 0, \\ X'(1) = 0. \end{cases}
\]

\subsection{ПЕРЕТВОРЕННЯ СКЛАДНИХ РІВНЯНЬ У ПРОСТУ
ФОРМУ}\label{ux43fux435ux440ux435ux442ux432ux43eux440ux435ux43dux43dux44f-ux441ux43aux43bux430ux434ux43dux438ux445-ux440ux456ux432ux43dux44fux43dux44c-ux443-ux43fux440ux43eux441ux442ux443-ux444ux43eux440ux43cux443}

МЕТА ЛЕКЦІЇ Показати, як PPP для функції \(u(x, t)\) можна перетворити
на простіший PPP для нової функції \(w(x, t)\). Вибір правильної
трансформації зазвичай базується на ненавмисному намірі. У нашій лекції

\[
u_t = \alpha^2 u_{xx} - \beta u,
u_t = \alpha^2 u_{xx} - v u_x, \quad (v = \text{const})
\]

переходимо до найпростішого рівняння теплопровідності

\[
w_1 = \alpha^2 w_{xx}
\]

внаслідок відповідних перетворень

\[
u(x, t) = e^{-\beta t} w(x, t)
u(x, t) = e^{\epsilon t (1 - \delta t - 2) 2\alpha^2} w(x, t).
\]

\begin{verbatim}
Розв'язавши просте рівняння теплопровідності за функцією \omega(x,t) , ми можемо знайти розв'язки початкових рівнянь, перетворивши функцію \omega(x,t) відповідно до цих формул (звісно, ці перетворення слід застосовувати і до початкових і граничних умов).
\end{verbatim}

Після прочитання двох останніх лекцій читач може створити враження, що
метод розділення змінних застосовується \emph{лише} до рівняння

\[
u_1 = \alpha^2 u_{xx}
\]

Насправді рівняння теплопотоку є найпростішим, але не єдиним
параболічним рівнянням, яке можна розв'язати методом розділення змінних.
Як уже згадувалося, якщо рівняння є лінійним і однорідним, його можна
розв'язати, поділивши змінні. Наприклад, двовимірна задача
теплопровідності в колі описується рівнянням

\[
u_t = \alpha^2 \left[ u_{rr} + \frac{1}{r} u_r + \frac{1}{r^2} u_{\theta\theta} \right],
\]

І хоча це рівняння з \emph{змінними коефіцієнтами}, його можна звести до
трьох звичайних диференціальних рівнянь.

У цій лекції ми покажемо читачеві, що іноді, замість того, щоб
безпосередньо розв'язувати диференціальне рівняння в похідних, краще
перетворити його на простіше рівняння, яке можна розв'язати методом
розділення змінних або іншим методом. Наступний приклад ілюструє цей
підхід.

\section{Перетворення проблеми з перенесенням бічної поверхні у проблему
з термічно ізольованою бічною
поверхнею}\label{ux43fux435ux440ux435ux442ux432ux43eux440ux435ux43dux43dux44f-ux43fux440ux43eux431ux43bux435ux43cux438-ux437-ux43fux435ux440ux435ux43dux435ux441ux435ux43dux43dux44fux43c-ux431ux456ux447ux43dux43eux457-ux43fux43eux432ux435ux440ux445ux43dux456-ux443-ux43fux440ux43eux431ux43bux435ux43cux443-ux437-ux442ux435ux440ux43cux456ux447ux43dux43e-ux456ux437ux43eux43bux44cux43eux432ux430ux43dux43eux44e-ux431ux456ux447ux43dux43eux44e-ux43fux43eux432ux435ux440ux445ux43dux435ux44e}

Розглянемо наступне завдання:

(8.1)

\[
u_{t} = \alpha^{2}u_{xx} - \beta u, \quad 0 < x < 1, \quad 0 < t < \infty,
\begin{cases} u(0, t) = 0, \\ u(1, t) = 0, \\ u(x, 0) = \varphi(x), \quad 0 \le x \le 1, \end{cases}
\]

(8.1)

\[
(FY) \quad
\begin{cases} u(0, t) = 0, \\ u(x, 0) = \varphi(x), \quad 0 \le x \le 1, \end{cases}
\]

\begin{verbatim}
де термін -\beta u описує потік тепла через бічну поверхню стрижня (рис. 8.1).
\end{verbatim}

\pandocbounded{\includegraphics[keepaspectratio]{images/_page_59_Figure_7.jpeg}}

\begin{verbatim}
Рис. 81. Задача теплопровідності, описана рівнянням u_t = \alpha^t u_{\lambda\lambda} - \beta u : \alpha — це втрата гепле через поверхню Боха; \delta – дифузія вздовж стрижня.
\end{verbatim}

\begin{verbatim}
Наша мета — ввести нову температурну w\left(x,\ t\right) замість u\left(x,\ t\right) , щоб рівняння для w\left(x,\ t\right) стало простішим за початкове рівняння
\end{verbatim}

\[
u_t = \alpha^2 u_{xx} - \beta u
\]

.

Це досить загальний метод розв'язання диференціальних рівнянь у
частинних похідних, але вибір відповідного перетворення базується на
інтуїтивних уявленнях про поведінку розв'язків початкового рівняння.
Наприклад, у задачі (8.1) температура \(u(x, t)\) у кожній точці \(x_0\)
змінюється внаслідок двох факторів:

\begin{verbatim}
- 1) дифузія тепла вздовж стрижня (через компонентний \alpha^2 u_{xx} ),
\end{verbatim}

\begin{itemize}
\tightlist
\item
  \begin{enumerate}
  \def\labelenumi{\arabic{enumi})}
  \setcounter{enumi}{1}
  \tightlist
  \item
    теплопередача через бічну поверхню стрижня (через член RI).
  \end{enumerate}
\end{itemize}

\begin{verbatim}
Найважливішим елементом нашого міркування є наступне: якщо не було дифузії вздовж стрижня ( \alpha = 0 ),
\end{verbatim}

тоді температура в кожній точці стрижня експоненційно знижується до нуля
згідно із законом

\[
u(x_0 t) = u(x_0, 0) e^{-\beta t}
\]

.

Виходячи з цього спостереження, запитаємо, чи можна розв'язок проблеми
(8.1) представити як утворення двох факторів

\[
(8.2) u(x, t) = e^{-\beta t}w(x, t)
\]

або:

\begin{verbatim}
\begin{pmatrix}
\end{verbatim}

Температура елемента з \textbackslash{}

Неізольовані боком

Поверхня

де \(w(x, t)\) описує розподіл температури, що виникає лише через
дифузію. Давайте подивимось, що станеться, якщо підставити цей вираз для
температури у задачу (8.1). Після виконання відповідних розрахунків
(читач може зробити їх сам), ми отримаємо

\begin{verbatim}
Нова задача майже точно така ж, як і оригінальна, але більше не має додатку 6u. Тому замість розв'язання задачі (8.1) ми розв'язуємо перетворену задачу (8.3), а потім множимо отриманий розв'язок w(x, t) на e^{-\beta t} . У нашому випадку розв'язок задачі (8.3) вже було отримано раніше методом розділення змінних, тому ми можемо одразу записати
\end{verbatim}

(8.4)

\[
w(x, t) = \sum_{n=1}^{\infty} a_n e^{-(n\pi\alpha)^{\frac{n}{2}}t} \sin(n\pi x),
a_n = 2 \int_0^{\infty} \phi(\xi) \sin(n\pi \xi) d\xi,
\]

і, отже, розв'язок початкової задачі (8.1) має вигляд

\[
u(x, t) = e^{-\beta t}w(x, t).
\]

Нижче, у розділі «Нотатки», наведено ще один приклад проблеми,
розв'язаної цим методом.

\section{ПРИМІТКИ}\label{ux43fux440ux438ux43cux456ux442ux43aux438}

\begin{enumerate}
\def\labelenumi{\arabic{enumi}.}
\tightlist
\item
  Вирішити проблему
\end{enumerate}

\[
\begin{aligned}
u_t &= u_{xx} - u, && 0 < x < 1,\ 0 < t < \infty, \\
u(0,t) &= 0,\quad u(1,t) = 0, && 0 < t < \infty, \\
u(x,0) &= \sin(\pi x) + 0.5\sin(3\pi x), && 0 \leqslant x \leqslant 1.
\end{aligned}
\]

За методом, описаним вище, ми

\begin{itemize}
\tightlist
\item
  \begin{enumerate}
  \def\labelenumi{\alph{enumi})}
  \tightlist
  \item
    тимчасово відкинути конвективний член u у початковому рівнянні.
  \end{enumerate}
\item
  \begin{enumerate}
  \def\labelenumi{\alph{enumi})}
  \setcounter{enumi}{1}
  \tightlist
  \item
    записати розв'язок змішаної задачі без конвективного члена
  \end{enumerate}
\end{itemize}

\[
u(x, t) = e^{-\pi^2 t} \sin(\pi x) + 0.5e^{-(3\pi)^2 t} \sin(3\pi x),
\]

\begin{enumerate}
\def\labelenumi{\alph{enumi})}
\setcounter{enumi}{2}
\tightlist
\item
  помножимо це розв'язання на конвективний фактор
  \(e^{-\beta t}=e^{-t}\) отримаємо остаточне розв'язання початкової
  задачі
\end{enumerate}

\[
u(x, t) = e^{-t} \left[ e^{-\pi^2 t} \sin(\pi x) + 0.5e^{-(3\pi)^2 t} \sin(3\pi x) \right].
\]

\begin{enumerate}
\def\labelenumi{\arabic{enumi}.}
\setcounter{enumi}{1}
\tightlist
\item
  Рівняння конвективної дифузії
\end{enumerate}

\[
u_t = \alpha^2 u_{xx} - v u_x
\]

(коефіцієнт і є константою) також можна перетворити у вигляд

\[
w_t = \alpha^2 w_{xx}.
\]

Відповідне перетворення має вигляд

\[
u(x, t) = e^{v(x-vt/2)/2\alpha^2}w(x, t).
\]

Це перетворення підкреслює експоненціальний фактор, який повністю
зумовлений рухом середовища. Зверніть увагу, що цей множник описує
хвилю, що поширюється зліва направо зі швидкістю \(\upsilon/2\) . Читач
матиме можливість скористатися цією трансформацією.

\section{ЗАВДАННЯ}\label{ux437ux430ux432ux434ux430ux43dux43dux44f-1}

\begin{enumerate}
\def\labelenumi{\arabic{enumi}.}
\tightlist
\item
  Розв'язати задачу конвективної дифузії
\end{enumerate}

\begin{verbatim}
\begin{aligned}
(\text{УЧП})\quad & u_t = u_{xx} - u_x, && 0 < x < 1,\ 0 < t < \infty, \\
(\text{ГУ})\quad & u(0,t) = 0,\quad u(1,t) = 0, && 0 < t < \infty, \\
(\text{НУ})\quad & u(x,0) = e^{x/2}, && 0 \leqslant x \leqslant 1.
\end{aligned}
\end{verbatim}

методом перетворення до простішого рівняння. На що схожий розв'язок?
Можна дати таку інтерпретацію цієї задачі: \(u(x,t)\) --- концентрація
деякої речовини в рухомому середовищі, яке рухається зліва направо зі
швидкістю \(v = 1\); початковий розподіл концентрації задається функцією
\(e^{x/2}\), а на межах концентрація підтримується нульовою. Чи
відповідає ваш розв'язок такій інтерпретації?

\begin{enumerate}
\def\labelenumi{\arabic{enumi}.}
\setcounter{enumi}{1}
\tightlist
\item
  Розв'язати задачу
\end{enumerate}

\begin{verbatim}
\begin{aligned}
(\text{УЧП})\quad & u_t = u_{xx} - u + x, && 0 < x < 1,\ 0 < t < \infty, \\
(\text{ГУ})\quad & u(0,t) = 0,\quad u(1,t) = 1, && 0 < t < \infty, \\
(\text{НУ})\quad & u(x,0) = 0, && 0 \leqslant x \leqslant 1.
\end{aligned}
\end{verbatim}

\begin{itemize}
\tightlist
\item
  а) перетворивши неоднорідні граничні умови на однорідні;
\item
  б) перетворивши вихідне рівняння на нове рівняння без члена \(-u\);
\item
  в) розв'язавши задачу, що виникає після цих перетворень.
\end{itemize}

\begin{enumerate}
\def\labelenumi{\arabic{enumi}.}
\setcounter{enumi}{2}
\tightlist
\item
  Розв'язати задачу
\end{enumerate}

\begin{verbatim}
\begin{aligned}
(\text{УЧП})\quad & u_t = u_{xx} - u, && 0 < x < 1,\ 0 < t < \infty, \\
(\text{ГУ})\quad & u(0,t) = 0,\quad u(1,t) = 0, && 0 < t < \infty, \\
(\text{НУ})\quad & u(x,0) = \sin(\pi x), && 0 \leqslant x \leqslant 1.
\end{aligned}
\end{verbatim}

застосовуючи метод розділення змінних безпосередньо до вихідного
рівняння, не виконуючи жодних попередніх перетворень. Чи збігається
отриманий розв'язок з розв'язком, знайденим раніше після заміни

\[
u(x,t) = e^{-t}w(x,t)
\]

?

\section{РОЗВ'ЯЗАННЯ НЕОДНОРІДНИХ УЧП МЕТОДОМ РОЗКЛАДУ ЗА ВЛАСНИМИ
ФУНКЦІЯМИ}\label{ux440ux43eux437ux432ux44fux437ux430ux43dux43dux44f-ux43dux435ux43eux434ux43dux43eux440ux456ux434ux43dux438ux445-ux443ux447ux43f-ux43cux435ux442ux43eux434ux43eux43c-ux440ux43eux437ux43aux43bux430ux434ux443-ux437ux430-ux432ux43bux430ux441ux43dux438ux43cux438-ux444ux443ux43dux43aux446ux456ux44fux43cux438}

ЦІЛЬ ЛЕКЦІЇ: Навчитися розв'язувати змішані задачі вигляду

\begin{verbatim}
\begin{aligned}
(\text{УЧП})\quad & u_t = \alpha^2 u_{xx} + f(x,t), && 0 < x < 1,\ 0 < t < \infty, \\
(\text{ГУ})\quad & \alpha_1 u_x(0,t) + \beta_1 u(0,t) = 0, && 0 < t < \infty, \\
& \alpha_2 u_x(1,t) + \beta_2 u(1,t) = 0, && 0 < t < \infty, \\
(\text{НУ})\quad & u(x,0) = \varphi(x), && 0 \leqslant x \leqslant 1.
\end{aligned}
\end{verbatim}

Розв'язок неоднорідного УЧП можна шукати у вигляді ряду

\[
u(x,t) = \sum_{n=1}^{\infty} T_n(t) X_n(x).
\]

де \(X_n(x)\) --- власні функції відповідної однорідної задачі

\begin{verbatim}
\begin{aligned}
(\text{УЧП})\quad & u_t = \alpha^2 u_{xx}, && 0 < x < 1,\ 0 < t < \infty, \\
(\text{ГУ})\quad & \alpha_1 u_x(0,t) + \beta_1 u(0,t) = 0, && 0 < t < \infty, \\
& \alpha_2 u_x(1,t) + \beta_2 u(1,t) = 0, && 0 < t < \infty, \\
(\text{НУ})\quad & u(x,0) = \varphi(x), && 0 \leqslant x \leqslant 1.
\end{aligned}
\end{verbatim}

А функції \(T_n(t)\) визначаються як розв'язки деяких звичайних
диференціальних рівнянь.

На лекції 6 ми ознайомилися з тим, як перетворювати неоднорідні граничні
умови на однорідні. На жаль, рівняння залишається гетерогенним, і ми
стоїмо перед завданням

\[
(9.1) \quad u_{t} = \alpha^{2} u_{xx} + f(x, t), \quad 0 < x < 1, \quad 0 < t < \infty,
(9.1) \quad (\Gamma Y) \quad
\begin{cases} \alpha_{1} u_{x}(0, t) + \beta_{1} u(0, t) = 0, \\ \alpha_{2} u_{x}(1, t) + \beta_{2} u(1, t) = 0, \\ u(x, 0) = \varphi(x), \quad 0 \le x \le 1. \end{cases}
\]

Мета нашої лекції --- опанувати метод розв'язання таких задач. Цей метод
аналогічний методу варіації довільних констант.

теорія звичайних диференціальних рівнянь називається методом розкладу за
власними функціями.

Ідея методу надзвичайно проста. Відомо, що розв'язок задачі (9.1) при
\(f(x, t)\) = 0 (так звана відповідна однорідна задача) записується як
ряд

\[
u\left(x,\ t\right)=\sum_{n=1}^{\infty}a_{n}e^{-\left(\lambda_{n}\alpha\right)^{2}}/X_{n}\left(x\right),
\]

де \(X_n(x)\) --- власні функції задачі Штурма-Ліувіля,

(9.2)

\[
X'' + \lambda^2 X = 0, \\ \alpha_1 X'(0) + \beta_1 X(0) = 0, \\ \alpha_2 X'(1) + \beta_2 X(1) = 0.
\]

Запитаймо себе, чи можна також шукати розв'язок гетерогенної проблеми
(9.1) у вигляді серії власних функцій, але більш загальної форми

\[
u(x, t) = \sum_{n=1}^{\infty} T_n(t) X_n(x).
\]

Існують певні фізичні підстави для такого припущення. Насправді терміни
\(f(x,t)\) описують джерело тепла, розташоване всередині стрижня. Це
означає, що залежність рішення від часу не буде виражена показниками

\[
e^{-(\lambda_n\alpha)^2 t}
\]

,

Це стосувалося однорідного завдання, але з деякими іншими функціями.

Щоб показати цей метод у дії і не ускладнювати подання деталями, ми
застосуємо його до розв'язання однієї дуже простої задачі.

\section{Построзв'язання методом демаркації за власними
функціями,}\label{ux43fux43eux441ux442ux440ux43eux437ux432ux44fux437ux430ux43dux43dux44f-ux43cux435ux442ux43eux434ux43eux43c-ux434ux435ux43cux430ux440ux43aux430ux446ux456ux457-ux437ux430-ux432ux43bux430ux441ux43dux438ux43cux438-ux444ux443ux43dux43aux446ux456ux44fux43cux438}

Розглянемо гетерогенну проблему

(9.3) (PY)

\[
\begin{aligned} u_t &= \alpha^2 u_{xx} + f(x, t), & 0 < x < 1, & 0 < t < \infty, \\ u(0, t) &= 0, & 0 < t < \infty, \\ u(1, t) &= 0, & 0 < t < \infty, \\ u(x, 0) &= \varphi(x), & 0 \leqslant x \leqslant 1. \end{aligned}
\]

Процес розв'язання задачі подається у вигляді послідовності наступних
кроків.

КРОК 1. Основна ідея методу полягає в розкладанні густини джерела
\(f(x, t)\) на ряд власними функціями

\[
f(x, t) = f_1(t) X_1(x) + f_2(t) X_2(x) + \dots + f_n(t) X_n(x) + \dots
\]

а також визначити реакції системи \(u_n\left(x,\,t\right)\) на вплив
кожного компонента \(f_n\left(t\right)X_n\left(x\right)\) . Якщо знайти
всі відповіді, рішення проблеми виглядатиме так:

\[
u(x, t) = \sum_{n=1}^{\infty} u_n(x, t).
\]

Головна складність цього методу --- розклад густини джерела на
компоненти \(f_n(t)\,X_n(x)\) . Виявляється, що \(X_n(x)\) факторами в
цій задачі є власними векторами системи Штурма--Ліувіля, які виникають
при розв'язанні відповідної однорідної задачі методом розділення
змінних, а саме

\[
\begin{array}{c} u_t = \alpha^2 u_{xx} \quad \text{(отметим, что } f(x, \ t) = 0), \\ u(0, \ t) = 0, \\ u(1, \ t) = 0, \\ u(x, \ 0) = \varphi(x). \end{array}
\]

Завдання Штурма-Ліувіля в цьому випадку виглядає так:

\[
X'' + \lambda^2 X = 0,
\]

\(X(0) = 0,\) \(X(1) = 0,\) а його розв'язки є функціями

\[
X_n(x) = \sin(n\pi x), \quad n = 1, 2, 3 \dots
\]

Отже, розклад на густину джерела представлений як ряд

\[
(9.5) \ f(x,t) = f_1(t) \sin(\pi x) + f_2(t) \sin(2\pi x) + \dots + f_n(t) \sin(n\pi x) + \dots
\]

\(\mathcal{U}\) нарешті, щоб знайти функції \(f_n(t)\) , помножте обидві
сторони останнього відношення на \(\sin(m\pi x)\) і інтегруйте від 0 до
1 за координатою x.

У результаті ми отримуємо

\[
\int_{0}^{1} f(x, t) \sin(m\pi x) dx = \sum_{n=1}^{\infty} f_{n}(t) \int_{0}^{1} \sin(m\pi x) \sin(n\pi x) dx =
= \frac{1}{2} f_{m}(t)
\]

або (після заміни t на p)

(9.6)

\[
f_{n}(t) = 2 \int_{0}^{1} f(x, t) \sin(n\pi x) dx.
\]

Останнє співвідношення встановлює співвідношення між коефіцієнтами
\(f_n\)(t) і часткою джерела тепла \(f(x, t)\).

КРОК 2. (Пошук \(u_n\)(x, t) = \(X_n\)(x) \(T_n\)(t) реакції на вхідний
вплив \(f_n\)(t) \(X_n\)(x) .)

Замінимо густину джерела на його розклад у рядок

\[
f(x, t) = \sum_{n=1}^{\infty} f_n(t) \sin(n\pi x)
\]

\begin{verbatim}
і намагатися знайти індивідуальні відповіді, тобто функції, T_{\mathbf{n}}(t) у
\end{verbatim}

\[
u(x, t) = \sum_{n=1}^{\infty} T_n(t) \sin(n\pi x).
\]

Оскільки розв'язок задачі має вигляд

\[
u(x, t) = \sum_{n=1}^{\infty} T_n(t) \sin(n\pi x),
\]

потім підставимо цей вираз у початкову задачу

\[
u_{t} = \alpha^{2} u_{xx} + \sum_{n=1}^{\infty} f_{n}(t) \sin(n\pi x),
u(0, t) = 0,
u(1, t) = 0,
u(x, 0) = \varphi(x),
\]

У результаті ми отримаємо

\[
\sum_{n=1}^{\infty} T'_n(t) \sin(n\pi x) = -\alpha^2 \sum_{n=1}^{\infty} (n\pi)^2 T_n(t) \sin(n\pi x) + + \sum_{n=1}^{\infty} f_n(t) \sin(\pi n\lambda),
\]

(9.7)

\[
\sum_{n=1}^{\infty} T_n(t) \sin 0 = 0
\]

(задовольняється ідентично),
\(\sum_{n=1}^{\infty} T_n(t) \sin (n\pi) = 0\) (задовольняється
ідентично), \(\sum_{n=1}^{\infty} T_n(0) \sin (n\pi x) = \phi(x)\) .

Переписування рівняння та початкової умови

(PPP)

\[
\sum_{n=1}^{\infty} [T'_n + (n\pi\alpha)^2 T_n - f_n(t)] \sin(n\pi x) = 0,
\]

(НУ)

\[
\sum_{n=1}^{\infty} T_n(0) \sin(n\pi x) = \varphi(x),
\]

можна побачити, що функції \(T_n(\ell)\) є розв'язками задачі Коші

(9.8)

\[
T_{n} + (n\pi\alpha)^{2} T_{n} = f_{n}(t),
T_{n}(0) = 2 \int_{0}^{1} \varphi(\xi) \sin(n\pi\xi) d\xi = a_{n}.
\]

Усі ці TAC легко розв'язати, а їхні рішення записуються у вигляді

\begin{enumerate}
\def\labelenumi{(\arabic{enumi})}
\setcounter{enumi}{98}
\tightlist
\item
\end{enumerate}

\[
T_n(t) = a_n e^{-(n\pi\alpha)^2 t} + \int_0^t e^{-(n\pi\alpha)^2 (t-\tau)} f_n(\tau) d\tau,
\]

З того місця, де легко отримати повне розв'язання проблеми (9.3):

(9.10)

\[
u(x, t) = \sum_{n=1}^{\infty} T_n(t) \sin(n\pi x) =
= \sum_{n=1}^{\infty} \left[ a_n e^{-(n\pi\alpha)^n t} \sin(n\pi x) \right] + \sum_{n=1}^{\infty} \left[ \sin(n\pi x) \int_0^t e^{-(n\pi\alpha)^n (t-\tau)} f_n(\tau) d\tau \right]
\]

Перехідна частина (обумовлена стаціонарною частиною (умовлена початковою
умовою) правою частиною \(f(x, t)\)).

Розв'язок (9.10) показує, що температурна реакція складається з двох
частин: перша частина зумовлена початковими умовами, а друга ---
густиною джерела \(f(x,t)\). Зверніть увагу, що термін «стаціонарна
частина» не є найкращим описом другої частини температурної реакції,
оскільки він не означає, що ця частина не залежить від часу
(стаціонарний стан може бути періодичним, якщо \(f(x, t)\) є періодичною
функцією часу t).

Проблема повністю вирішена. Перш ніж перейти до наступної теми, давайте
покажемо, як використовувати описаний вище метод для розв'язання

Конкретне завдання.

70

\section{Розв'язання конкретної задачі методом декомпозиції за власними
функціями}\label{ux440ux43eux437ux432ux44fux437ux430ux43dux43dux44f-ux43aux43eux43dux43aux440ux435ux442ux43dux43eux457-ux437ux430ux434ux430ux447ux456-ux43cux435ux442ux43eux434ux43eux43c-ux434ux435ux43aux43eux43cux43fux43eux437ux438ux446ux456ux457-ux437ux430-ux432ux43bux430ux441ux43dux438ux43cux438-ux444ux443ux43dux43aux446ux456ux44fux43cux438}

Розглянемо просту задачу

(9.11)

\[
\begin{aligned} u_t &= \alpha^2 u_{xx} + \sin{(3\pi x)}, & 0 < x < 1, \\ (\mathbf{FY}) & \begin{cases} u(0, t) &= 0, \\ u(1, t) &= 0, \end{cases} & 0 < t < \infty, \\ (\mathbf{HY}) & u(x, 0) &= \sin{(\pi x)}, & 0 \leq x \leq 1. \end{aligned}
\]

Наша мета --- обчислити коефіцієнти \(T_n(t)\) при розкладі

\[
u\left(x,\,t\right)=\sum_{n=1}^{\infty}\,T_{n}\left(t\right)\sin\left(n\pi x\right)
\]

(Власні функції цієї задачі \(X_n(x)\) збігатися з власними функціями
попереднього). Якщо підставити цю декомпозицію на початкову задачу,
отримаємо наступну задачу Коші для визначення функцій \(T_n(t)\) :

\[
T_n + (n\pi\alpha)^2 T_n = \int_n (t) = 2 \int_0^1 \sin(3\pi x) \sin(n\pi x) dx =
\begin{cases} 1 &n = 3, \\ 0 & n \neq 3, \end{cases}
T_n(0) = 2 \int_0^1 \sin(n\xi) \sin(n\pi\xi) d\xi =
\begin{cases} 1 &n = 1, \\ 0 & n \neq 1. \end{cases}
\]

Перепишем эти уравнения последовательно для n=1, 2\ldots{}

\begin{verbatim}
\begin{aligned}
(n=1)\quad & T_1' + (\pi\alpha)^2 T_1 = 0, && T_1(0) = 1, && T_1(t) = e^{-(\pi\alpha)^2 t}, \\
(n=2)\quad & T_2' + (2\pi\alpha)^2 T_2 = 0, && T_2(0) = 0, && T_2(t) = 0, \\
(n=3)\quad & T_3' + (3\pi\alpha)^2 T_3 = 4, && T_3(0) = 0, && T_3(t) = \frac{1}{(3\pi\alpha)^2}\left[1 - e^{-(3\pi\alpha)^2 t}\right], \\
(n\geqslant 4)\quad & T_n' + (n\pi\alpha)^2 T_n = 0, && T_n(0) = 0, && T_n(t) = 0.
\end{aligned}
\end{verbatim}

Отже, розв'язок задачі можна записати у вигляді

\[
u(x, t) = e^{-(\pi\alpha)^2 t} \sin(\pi x) + \frac{1}{(3\pi\alpha)^2} [1 - e^{-(3\pi\alpha)^2 t}] \sin(3\pi x).
\]

(9.12)

\section{ЗАУВАЖЕННЯ}\label{ux437ux430ux443ux432ux430ux436ux435ux43dux43dux44f-1}

\begin{enumerate}
\def\labelenumi{\arabic{enumi}.}
\item
  Метод розкладу власних функцій є одним із найпотужніших методів
  розв'язання неоднорідних диференціальних рівнянь у похідних похідних.
  Пізніше, коли ми ознайомимося з інтегральними перетвореннями,
  побачимо, що існують й інші методи розв'язання задач такого типу.
\item
  Власні функції \(X_n\)(x) , за допомогою якої виконується розклад,
  залежать від задачі, яку потрібно розв'язати, точніше від рівнянь і
  граничних умов. Читача просять розв'язати задачу 4 наприкінці лекції,
  щоб перевірити, чи дійсно він навчився знаходити власні функції
  \(X_n\)(x) .
\item
  Якщо читач пам'ятає основні факти з теорії звичайних диференціальних
  рівнянь, то він знає, що розв'язки, що відповідають неоднорідним
  членам вигляду
\end{enumerate}

\[
P_n(x) e^{\alpha x}
\begin{cases} \sin{(\beta x)}, \\ \cos{(\beta x)}, \end{cases}
\]

можно найти методом неопределенных коэффициентов. То же самое и здесь.
Задача (9.11) может быть решена таким же способом.

\section{ЗАДАЧИ}\label{ux437ux430ux434ux430ux447ux438}

\begin{enumerate}
\def\labelenumi{\arabic{enumi}.}
\item
  Подставьте разложение (46.8) в задачу (46.7) и получите
  последовательность задач \(P_0,\ P_1,\ P_2,\ \dots\).
\item
  Покажите, что нелинейную задачу
\end{enumerate}

\[
\begin{aligned}
\Delta u + u^2 &= 0, && 0 \leqslant r < 1,\ 0 \leqslant \theta < 2\pi, \\
u(1,\theta) &= \cos\theta, && 0 \leqslant \theta < 2\pi,
\end{aligned}
\]

можно свести к последовательности линейных задач \(P_0,\ P_1,\ P_2\).

\begin{enumerate}
\def\labelenumi{\arabic{enumi}.}
\setcounter{enumi}{2}
\item
  Подставьте приближенное решение (46.5) в нелинейную задачу 2 и оцените
  его точность.
\item
  Решите задачу \(P_1\) из примера с возмущенной границей и проверьте,
  насколько хорошо приближенное решение
\end{enumerate}

\[
u(r,\theta) = u_0(r,\theta) + \frac{1}{4}u_1(r,\theta)
\]

удовлетворяет соотношениям

\[
\Delta u = 0,
u\left(1 + \frac{1}{4}\sin\theta,\ \theta\right) = \cos\theta.
\]

\section{Лекція 47}\label{ux43bux435ux43aux446ux456ux44f-47}

\section{РОЗВ'ЯЗАННЯ ДИФЕРЕНЦІАЛЬНИХ РІВНЯНЬ У ПОХІДНИХ ПОХІДНИХ ЗА
ДОПОМОГОЮ МЕТОДУ КОНФОРМНИХ
ВІДОБРАЖЕНЬ}\label{ux440ux43eux437ux432ux44fux437ux430ux43dux43dux44f-ux434ux438ux444ux435ux440ux435ux43dux446ux456ux430ux43bux44cux43dux438ux445-ux440ux456ux432ux43dux44fux43dux44c-ux443-ux43fux43eux445ux456ux434ux43dux438ux445-ux43fux43eux445ux456ux434ux43dux438ux445-ux437ux430-ux434ux43eux43fux43eux43cux43eux433ux43eux44e-ux43cux435ux442ux43eux434ux443-ux43aux43eux43dux444ux43eux440ux43cux43dux438ux445-ux432ux456ux434ux43eux431ux440ux430ux436ux435ux43dux44c}

МЕТА ЛЕКЦІЇ: Показати, як деякі двовимірні граничні задачі можна
перетворити на інші, простіші задачі за допомогою конформного
відображення. Припустимо, наприклад, що ми хочемо розв'язати рівняння
Лапласа

Щоб візуально зрозуміти, як працює ця функція, ми радимо читачеві
побудувати її за кількома лініями y=c.

У другому прикладі ми перетворюємо область між двома неконцентричними
колами у кільце.

Задача Діріхле в області між двома некопоцентричними колами

Припустимо, ви хочете знайти потенціал між двома колами

з певною граничною умовою в області комплексної форми (площина змінних x
і y). Цю задачу можна перетворити на нову, в якій потрібно знайти
розв'язок рівняння Лапласа

.

Тут z --- це комплексна змінна (з деякої області площини z), а w ---
нова комплексна змінна, яка отримає значення відповідно до формули w =
\(f(z)\) у новій комплексній площині w.

Наприклад, комплексна функція, визначена в \emph{першому квадранті}
комплексної z-площини, \(w=z^2\) відображає цю область у всю півплощину
v\textgreater0 у комплексній площині w.

Таблиця 47 2

{\def\LTcaptype{none} % do not increment counter
\begin{longtable}[]{@{}
  >{\raggedright\arraybackslash}p{(\linewidth - 4\tabcolsep) * \real{0.2342}}
  >{\raggedright\arraybackslash}p{(\linewidth - 4\tabcolsep) * \real{0.2252}}
  >{\raggedright\arraybackslash}p{(\linewidth - 4\tabcolsep) * \real{0.5405}}@{}}
\toprule\noalign{}
\begin{minipage}[b]{\linewidth}\raggedright
Площина d
\end{minipage} & \begin{minipage}[b]{\linewidth}\raggedright
\(z \longrightarrow z^2\)
\end{minipage} & \begin{minipage}[b]{\linewidth}\raggedright
Площина S
\end{minipage} \\
\midrule\noalign{}
\endhead
\bottomrule\noalign{}
\endlastfoot
0 & & 0 \\
t1 → t & & -1(2) \\
Вісь x Перший квадрант & & Позитивна дійсна піввісь верхньої
півплощини \\
\end{longtable}
}

Таблиця 47.2 показує відповідність, яка встановлюється шляхом
відображення \(w=z^2\) між деякими точками на площинах z і w.

\begin{verbatim}
Рис. 47.1. Площина комплексної змінної  \boldsymbol{z}  та некольорових формул бігунів Р И С У Н
\end{verbatim}

\[
\begin{aligned} & (\text{УЧП}) & \phi_{xx} + \phi_{yy} = 0, & \text{внутри } D, \\ & (\text{ГУ}) & \begin{cases} \varphi(x, y) = 1 & \text{на } x^2 + y^2 = 1, \\ \varphi(x, y) = 2 & \text{на } (x - 1)^2 + y^2 = 9. \end{cases} \end{aligned}
\]

\begin{verbatim}
ОК 47.2. Відображення першого квадранта площини  {\bf z}  у верхню півплощину  {\bf w}  .
\end{verbatim}

Інакше кажучи, якщо ми говоримо про комплексісні функції змінної w =
\(f(z)\), то можна розглянути питання, як криві площини z перетворюються
на криві в площині

Це питання не ставиться в теорії функцій дійсної змінної, з якою читач
уже добре знайомий.

Перш ніж рухатися далі, давайте розглянемо, як w=z\^{}2 складне
відображення типів у еквівалентній дійсній формі. Дійсна форма показує,
як початкові координати (x,y) у площині z перетворюються на нові
координати (u,v) у площині w w=z\^{}2 .

(HY)

І давайте окремо прирівняємо реальні та уявні частини в цьому
співвідношенні. У результаті ми отримуємо

\[
w = 2t \left[ \frac{z-s}{z-t} \right],
\]

(действительная форма отображения w=z\^{}2 ).

Ця форма нам знадобиться пізніше. У цьому прикладі це цікаво, оскільки
показує, як гіперболи в площині z переходять до координатних прямих u =
const і v = const у площині w.

Тепер перейдемо до вивчення одного конкретного типу функцій комплексної
змінної, які виконують конформні відображення.

\section{Определение конформного
отображения}\label{ux43eux43fux440ux435ux434ux435ux43bux435ux43dux438ux435-ux43aux43eux43dux444ux43eux440ux43cux43dux43eux433ux43e-ux43eux442ux43eux431ux440ux430ux436ux435ux43dux438ux44f}

\begin{verbatim}
Відображення w=f(z) комплексної площини z у комплексну площину w називається конформним у точці  z_{\bullet}  площини z, якщо похідна  f'(z_{\bullet}) \neq 0  . Відображення f(z) називається конформним у регіоні D, якщо  f'(z) \neq 0  в кожній точці регіону D.
\end{verbatim}

\begin{verbatim}
Наприклад, відображення  f(z) = z^2  конформно всюди, крім z = 0, оскільки  z'(z) = 2z \neq 0  для всіх  z \neq 0  . З іншого боку, відображення  e^z  конформно по всій площині z, оскільки скрізь  f'(z) = e^z \neq 0  . Яке призначення конформних відображень? Відповідь полягає в тому, що при розв'язанні рівняння Лапласа  \phi_{xx} + \phi_{yy} = 0  в межах деякої області площини змінних x і y, ми можемо розглядати цю площину як площину комплексної змінної z. Розглянемо комплексне відображення w = f(z) площини z у площину w, у якій введені координати (u, v). Коли w = f(z  \phi_{xy} + \phi_{uy} = 0  ) відображається, рівняння Лапласа  \varphi_{xx} + \varphi_{yy} = 0  перетворюється у нове диференціальне рівняння в похідних похідних, яке залежить від нових координат і v. To — рівняння, яке  \phi_{xx} + \phi_{uu} = 0  повернеться у рівняння  \varphi_{nn} + \varphi_{nn} = 0  . Таким чином, ми придумаємо ідею знайти конформне відображення, яке переносить область з комплексною межею у область з простою межею (нагадаємо, що відображення  w=z^*  переносить межу першого квадранта площини z у фактичну вісь площини w).
\end{verbatim}

Ось кілька прикладів, які покажуть, як можна розв'язати диференціальні
рівняння у похідних за допомогою конформних відображень.

Уравнение Лапласа в верхней полуплоскости.

Припустимо, що нам потрібно розв'язати наступну задачу Діріхле у верхній
половині видимості (рис. 47.3):

(47.1)

Один із способів розв'язання цієї проблеми --- застосувати перетворення
Фур'є до змінної x, але набагато краще конформно відобразити верхню
половину площини у площину (u, v). Читачеві не буде складно побачити, що
конформне відображення

\[
u = \gamma [(x-s)(x-t)-\gamma^2],
\]

Вона має такі властивості (див. RNS. 47.4).

\begin{enumerate}
\def\labelenumi{\arabic{enumi}.}
\tightlist
\item
  Вона відображає верхню півплощину площини z у смуги
  \(-\infty < u < \infty\) , \(0 < v < \pi\) , у площині w.
\end{enumerate}

РИСУНОК 47.3 Проблема граничного значення у верхній половині площі.

\begin{verbatim}
- 2. Ребро прямої y=0, -1 < x < 1, у площині z (де потенціал дорівнює 1), переходить у пряму  v=\pi,  -  \infty < u < \infty  , у площині w. - 3. Два промені x>1, y=0, і x<-1, y=0, у площині z проходять відповідно у позитивний і негативний
\end{verbatim}

Читач повинен бути впевнений, що може підтвердити ці твердження.

\[
\quad 0 \le x \le L,
\]

Рис. 47.4. Перетворення складної задачі на легку за допомогою
конформного відображення.

Важливість цих властивостей полягає в тому, що вони перетворюють
початкову вадаху (47.1) на дуже просту задачу

(47.2)

\[
u(x, t) = X(x) T(t).
\]

Виявляється, що в цьому випадку хвилі, що рухаються, відбиваються від
меж так, що отримані коливання стають непоширюючими, а зберігають форму
в одному положенні, тобто Р

\[
T'' - \alpha^2 \lambda T = 0,
\]

ИС. 20.1. Три типові стоячі хвилі: \(X(x)\) \(T(t)\).

Декілька прикладів стоячих хвиль показано на рис. 20.1. Якщо профілі цих
стоячих хвиль відомі \(X_n\)(x) і природа коливань кожної хвилі
\(T_n\)(t) , то розв'язок проблеми коливань гітарних струн є
суперпозицією найпростіших коливань \(X_n\)(x),\(T_n\)(t) і

сывается в виде

\[
X'' - \lambda X = 0,
\]

\begin{verbatim}
Залишилося лише обрати шанси  c_n  так, щоб при t = 0 початкові умови  u\left(x,\ 0\right)=f\left(x\right)  і  u_t\left(x,\ 0\right)=g\left(x\right)  були виконані.
\end{verbatim}

Задачу о колебаниях гитарной струны мы решим методом разделения
переменных.

\section{Розв'язання проблеми коливань обмеженої струни методом
розділення
змінних}\label{ux440ux43eux437ux432ux44fux437ux430ux43dux43dux44f-ux43fux440ux43eux431ux43bux435ux43cux438-ux43aux43eux43bux438ux432ux430ux43dux44c-ux43eux431ux43cux435ux436ux435ux43dux43eux457-ux441ux442ux440ux443ux43dux438-ux43cux435ux442ux43eux434ux43eux43c-ux440ux43eux437ux434ux456ux43bux435ux43dux43dux44f-ux437ux43cux456ux43dux43dux438ux445}

Чтобы решить смешанную задачу

(УЧП)

\[
u(x, t) = [C \sin(\beta x) + D \cos(\beta x)][A \sin(\alpha \beta t) + B \cos(\alpha \beta t)]
\beta_n = \frac{n\pi}{L}, \quad n = 0, 1, 2, \ldots
u_n(x, t) = X_n(x) T_n(t) = \sin(n\pi x/L) [a_n \sin(n\pi at/L) + b_n \cos(n\pi at/L)],
u_n(x, t) = R_n \sin\left(\frac{n\pi x}{L}\right) \cos\left(\frac{n\pi\alpha (t - \delta_n)}{L}\right).
\begin{cases} u(0, t) = 0, \\ u(L, t) = 0, \end{cases}
u(x, t) = \sum_{n=1}^{\infty} \sin(n\pi x/L) \left[ a_n \sin(n\pi \alpha t/L) + b_n \cos(n\pi \alpha t/L) \right]
H(x-a) =
\]

якщо потенціал на внутрішньому колі дорівнює одинці, а на зовнішньому
колі дорівнює двом (рис. 47.5). Іншими словами, нам потрібно розв'язати
проблему

(47.3)

\[
u(x, \theta) = f(x),
\]

де s=-0,146, а t=-6,85. У своїй еквівалентній дійсній формі це буде

(47.5)

\[
\sum_{n=1}^{\infty} b_n \sin(n\pi x/L) = f(x),
\begin{cases} 0, & x < a, \\ 1, & x \geqslant a, \end{cases}
\sum_{n=1}^{\infty} a_n (n\pi \alpha/L) \sin(n\pi x/L) = g(x).
\]

найдем сначала стоячне волны, т. е. решения вида

\[
f(x) = \mathcal{F}^{-1}[F] = \frac{1}{\sqrt{2\pi}} \int_{-\infty}^{+\infty} F(\omega) e^{-i\omega x} d\omega \qquad F(\omega) = \mathcal{F}[f] = \frac{1}{\sqrt{2\pi}} \int_{-\infty}^{+\infty} f(x) e^{-i\omega x} dx
\]

Підставляючи цей вираз у рівняння (20.1) і ділячи змінні, отримуємо дві
звичайні диференціальні рівняння

\[
\int_{0}^{L} \sin(m\pi x/L) \sin(n\pi x/L) dx =
\begin{cases} 1, |x| < a \\ 0, |x| > a \end{cases}
\begin{cases} 1, |x| < 1 \\ 0, |x| > 1 \end{cases}
a_{n} = \frac{2}{n\pi\alpha} \int_{0}^{L} g(x) \sin(n\pi x/L) dx, \\ b_{n} = \frac{2}{L} \int_{0}^{L} f(x) \sin(n\pi x/L) dx.
(1+x^2)^{-1}
u(x, t) = \sum_{n=1}^{\infty} \sin(n\pi x/L) \left[ a_n \sin(n\pi \alpha t/L) + b_n \cos(n\pi \alpha t/L) \right],
xe^{-a+x^4}, a > 0
\]

\begin{verbatim}
\begin{aligned} (\text{Y} \forall \Pi) & \qquad \phi_{\text{vx}} + \phi_{yy} = 0, \quad 0 < x < \infty, \quad 0 < y < \infty, \\ \left\{ \begin{array}{l} \phi(x, \ 0) = \begin{cases} 1, & 0 < x < 1, \\ 0, & 1 \leqslant x < \infty, \end{cases} \\ \phi(0, \ y) = \begin{cases} 1, & 0 < y < 1, \\ 0, & 1 \leqslant y < \infty. \end{cases} \end{aligned} \end{array}
\end{verbatim}

\[
u(x, t) = \sum_{n=1}^{\infty} b_n \sin(n\pi x/L) \cos(n\pi\alpha t/L)
u(x, \theta) = f(x)
\]

Зверніть увагу на зображення.

\[
\frac{a}{x^2+a^2}
f(x) = \sum_{n=1}^{\infty} b_n \sin(n\pi x/L).
\]

\begin{enumerate}
\def\labelenumi{\arabic{enumi}.}
\setcounter{enumi}{3}
\tightlist
\item
  Розв'язати змішану задачу Діріхле-Неймана з кутовим порошком за
  допомогою розчину під кутом 45°
\end{enumerate}

\[
\frac{2ax}{(x^2+a^2)^2}
\]

\begin{verbatim}
ПРИМІТКА. Комплексна функція  w = \ln z = \ln |z| + i \arg(z)  відображає промінь  \theta = c_1  у площині z у пряму  v = c_1  , а коло  r = c_2  у пряму  u = \ln c_2  у площині w (див. рисунки).
\end{verbatim}

\[
u_n(x, t) = b_n \sin(n\pi x/L) \cos(n\pi\alpha t/L).
\]

\section{ТАБЛИЦІ ІНТЕГРАЛЬНИХ
ПЕРЕТВОРЕНЬ}\label{ux442ux430ux431ux43bux438ux446ux456-ux456ux43dux442ux435ux433ux440ux430ux43bux44cux43dux438ux445-ux43fux435ux440ux435ux442ux432ux43eux440ux435ux43dux44c}

Таблиця A: Експоненціальне перетворення Фур'є,

Таблиця B: Перетворення синуса Фур'є.

Таблиця St Cosine Фур'є-перетворення.

Таблиця Dr скінченне перетворення Фур'є, таблиця та скінченне косинусне
перетворення Фур'є.

Таблиця D: Трансформація Лапласа.

\begin{verbatim}
# Прийняті позначення  \delta\left(x\right)  — дельта-функція,
\end{verbatim}

\[
\begin{cases} \cos(ax), |x| > \pi/2a \\ 0, |x| > \pi/2a \end{cases}
u(x, t) = \sin(\pi x/L)\cos(\pi \alpha t/L) + 0.5\sin(3\pi x/L)\cos(3\pi \alpha t/L) + 0.25\sin(5\pi x/L)\cos(5\pi \alpha t/L).
\begin{cases} 1-|x|, |x| < 1 \\ 0, |x| > 1 \end{cases}
R_n \sin(n\pi x/L) \cos[n\pi\alpha(t-\delta_n)/L],
\omega_n = \frac{n\pi\alpha}{L} = \frac{n\pi}{L} \sqrt{\frac{T}{\rho}},
\]

Воспользовавшись соотношениями ортогональности

\[
\frac{2}{\pi} \frac{\sin a\omega}{\omega}
\]

\begin{enumerate}
\def\labelenumi{\arabic{enumi}.}
\setcounter{enumi}{8}
\tightlist
\item
\end{enumerate}

\[
u(x, 0) = \sin(\pi x/L) + 0.5 \sin(3\pi x/L),
\]

\begin{enumerate}
\def\labelenumi{\arabic{enumi}.}
\setcounter{enumi}{9}
\tightlist
\item
\end{enumerate}

\[
\frac{2}{\pi} \frac{\sin a\omega}{(\omega^2+a^2)^2}
\]

находим коэффициенты \(a_n\) и \(b_n\) (20.5)

\[
u(x, 0) = 0,
\]

Итак, задача решена. Ее решение записывается в виде

(20.6)

\(\frac{\pi}{2}\)e\^{}\{-a+\omega\}\[
\]

а шанси \$ a\_n \$ і \$ b\_n \$ визначаються за формулами (20,5). Перед
завершенням лекції дозвольте зробити кілька зауважень.

\end{document}
